\documentclass[12pt, a4paper]{article}
\usepackage[utf8]{inputenc}
\usepackage{geometry}
\usepackage{listings}
\usepackage{xcolor}
\usepackage{graphicx}
\usepackage{tabularx}
\usepackage{float}
\usepackage{hyperref}

% Page Margins
\geometry{top=1in, bottom=1in, left=1in, right=1in}

% Styling
\definecolor{codegreen}{rgb}{0,0.6,0}
\definecolor{codegray}{rgb}{0.5,0.5,0.5}
\definecolor{codepurple}{rgb}{0.58,0,0.82}
\definecolor{backcolour}{rgb}{0.95,0.95,0.92}

\lstdefinestyle{style}{
    backgroundcolor=\color{backcolour},   
    commentstyle=\color{codegreen},
    keywordstyle=\color{magenta},
    numberstyle=\tiny\color{codegray},
    stringstyle=\color{codepurple},
    basicstyle=\ttfamily\footnotesize,
    breakatwhitespace=false,         
    breaklines=true,                 
    captionpos=b,                    
    keepspaces=true,                 
    numbers=left,                    
    numbersep=5pt,                  
    showspaces=false,                
    showstringspaces=false,
    showtabs=false,                  
    tabsize=2
}

\lstset{style=style}

\title{\textbf{CS327: Compilers} \\ Assignment 2: Lexical Analyzer}
\author{
    \textbf{Group Members:} \\
    Gaurav Budhwani (22110085) \\
    John Twipraham Debbarma (22110108) \\
    Harshit (22110095)
}
\date{\today}

\begin{document}

\maketitle
\tableofcontents
\newpage

\section{Introduction}
This report documents the implementation of the Lexical Analyzer for the \textbf{OPEX} programming language. The lexer uses \textbf{Flex} to tokenize the source code based on the language specifications.

\section{Regex Patterns}
The following regular expressions are used to define the tokens:

\begin{itemize}
    \item \textbf{Identifier:} \texttt{[a-zA-Z\_][a-zA-Z0-9\_]*}
    \item \textbf{Integer:} \texttt{[0-9]+}
    \item \textbf{Float:} \texttt{[0-9]+\textbackslash.[0-9]+([eE][+-]?[0-9]+)?}
    \item \textbf{String:} \texttt{"[\textasciicircum"\textbackslash n]*"}
    \item \textbf{Operators:} \texttt{:=, .., ++, --, +, -, *, /, \%, ==, !=, <, >, <=, >=, |>, **}
\end{itemize}

\section{Source Code}
\subsection{lexer.l}
The main Flex specification file handles token recognition.

\lstinputlisting[language=C, caption=lexer.l]{src/lexer.l}

\subsection{Makefile}
The build system compiles the lexer using \texttt{flex} and \texttt{gcc}.

\lstinputlisting[language=make, caption=Makefile]{Makefile}

\section{Test Cases}
We tested the lexer with various OPEX programs covering different language features.

\subsection{Test 1: Matrix Operations (sample.opx)}
This test checks matrix declaration, loops, and slicing.

\lstinputlisting[language=C, caption=sample.opx]{tests/sample.opx}

\textbf{Output (Snippet):}
\begin{verbatim}
KW_PROG: PROGRAM
ID: MatrixMagic
KW_COMPOSITE: STRUCT
ID: Pixel
...
\end{verbatim}

\subsection{Test 2: Prime Check (prime.opx)}
This test verifies boolean logic, conditionals, and arithmetic operators.

\lstinputlisting[language=C, caption=prime.opx]{tests/prime.opx}

\subsection{Test 3: Edge Detection (edge.opx)}
This test verifies kernel definitions and pipeline operators.

\lstinputlisting[language=C, caption=edge.opx]{tests/edge_detection.opx}

\subsection{Test 4: Palindrome Check (palindrome.opx)}
This test verifies relational operators (==, !=) and while loops.

\lstinputlisting[language=C, caption=palindrome.opx]{tests/palindrome.opx}

\subsection{Test 5: Matrix Multiplication (matrix\_mult.opx)}
This test demonstrates nested loops and 2D array access.

\lstinputlisting[language=C, caption=matrix\_mult.opx]{tests/matrix_mult.opx}

\subsection{Test 6: Recursive Fibonacci (fibonacci.opx)}
This test demonstrates global function definition and recursion.

\lstinputlisting[language=C, caption=fibonacci.opx]{tests/fibonacci.opx}

\end{document}
